% p3-02-notation

\section{P3 §2. Notation and Representation Hierarchy}

2. Notation and Representation Hierarchy

         Figure (fig-p3-ch2-notation). Notation triptych --- Constants / Sanctioned seeds / Lucas ring
                                                    lattice.

  2.1 Fundamental Constants and Parameters
  Throughout this paper the following constants are fixed:

  - $\varphi$ = (1 + 5)/(2) approx 1.61803398874989, the golden ratio;
  - pi approx 3.14159265358979, the ratio of a circle's circumference to its diameter;
  - e approx 2.71828182845905, the base of the natural logarithm;
  - psi = $\varphi$-1 = $\varphi$ - 1 approx 0.61803398874989, the reciprocal golden ratio.

  Note the fundamental identity $\varphi$2 = $\varphi$ + 1, which implies $\varphi$-1 = $\varphi$ - 1 and more generally
  the recurrence $\varphi$n = Fn $\varphi$ + Fn-1 for Fibonacci numbers Fn. The identity

  $\varphi$2 + $\varphi$-2 = 3 star

  is the Trinity anchor identity. It follows directly from $\varphi$2 = $\varphi$ + 1 and $\varphi$-2 = 2 - $\varphi$:

  $\varphi$2 + $\varphi$-2 = ($\varphi$ + 1) + (2 - $\varphi$) = 3.

  This is not a coincidence or a notational choice: it is the algebraic root of the
  multiplicative Trinity monomial grammar (Section 3). The factor 3 in n $\cdot$ 3k reflects
  precisely this sum, making the Trinity monomial system grounded in Z[$\varphi$]-arithmetic
  rather than arbitrary. Notably, the same golden ratio governs the structure of
  quasicrystalline materials --- in Shechtman et al.'s discovery, the linear scale between
  pentagonal diffraction features is exactly tau = $\varphi$ --- suggesting that the algebraic
  centrality of $\varphi$ in the Trinity framework resonates with its role in aperiodic physical
  order.

  The proof of (star) has Qed status in the t27 Coq proof base (C-01 in the Coq citation
  map of the PhD program's Appendix F; see Appendix E of this paper).

  2.2 Sanctioned Seeds and Forbidden Seeds
  The PhD program specifies sanctioned Fibonacci/Lucas initial seeds for experiments
  and formal computations:

  Sanctioned seeds: F17 = 1597, F18 = 2584, F19 = 4181, F20 = 6765, F21 = 10946, L7
  = 29, L8 = 47.

  Forbidden STROBE seeds: {42, 43, 44, 45}. Any numerical experiment in the
  computational program of Section 9 that uses these seeds as RNG initialisation values
  is deemed non-reproducible within the PhD framework and must be re-run with
  sanctioned alternatives.

  The Pellis expansion algorithm (Algorithm 4.1) applied to these seeds produces
  canonical coefficient sequences whose study is a primary focus of the computational
  program. Each sanctioned seed lies in Z[$\varphi$] (as an integer), so by Corollary 4.5 its Pellis
  expansion terminates exactly.

  2.3 The Pellis Hierarchy: Levels and Depths

        Definition 2.1 (Pellis depth-D expansion). A Pellis expansion of depth D of a
        real number x > 0 is a formal sum

  x = sumk=0D ck(D) $\varphi$-k + RD(x),

  where the coefficients ck(D) $\in$ Z$\geq$ 0 are determined by a hierarchical greedy algorithm
  (Algorithm 4.1), and RD(x) is the depth-D remainder satisfying 0 $\leq$ RD(x) < $\varphi$-(D+1).

  The hierarchical interference structure refers to the property that coefficients at depth
  d may encode cancellations between terms at depths d-1 and d+1, analogous to detail
  coefficients in a wavelet expansion. This is made precise in Section 4.2.

  2.4 The Trinity Monomial Lattice

        Definition 2.2 (Trinity monomial). A Trinity monomial is an expression of the
        form

  T(n,k,p,m,q) = n $\cdot$ 3k $\cdot$ $\varphi$p $\cdot$ pim $\cdot$ eq,

  where n $\in$ Z$\geq$ 1, and (k, p, m, q) $\in$ Z4. The quintuple lambda = (n, k, p, m, q) is called
  the Trinity label of the monomial.

  The set of all Trinity monomials

  T = { T(n,k,p,m,q) : n $\geq$ 1, (k,p,m,q) $\in$ Z4 }

  is conjecturally dense in (0, $\in$fty) under the working independence hypothesis (see
  Section 3.4). The factor 3 in the exponent 3k is not arbitrary: it is motivated by the
  Trinity anchor identity (star). Within Z[$\varphi$], the number 3 arises as the sum of $\varphi$2 and its
  conjugate $\varphi$-2; it is therefore the simplest non-unit, non-power-of-$\varphi$ element of Z[$\varphi$].

  2.5 The Lucas Ring Z[$\varphi$] as Discrete State Space

        Definition 2.3 (Lucas ring). The Lucas ring is Z[$\varphi$] = { a + b$\varphi$ : a, b $\in$ Z }, the
        ring of integers of the quadratic field Q(5). Lucas numbers Ln and Fibonacci
        numbers Fn both live in Z[$\varphi$] via

  Ln = $\varphi$n + psin, Fn = ($\varphi$n - psin)/(5),

  where psi = $\varphi$-1 = $\varphi$ - 1 is the Galois conjugate of $\varphi$. This is the standard Binet formula;
  the Fibonacci recurrence and the identity Ln = Fn-1 + Fn+1 are both consequences.

  The Lucas ring plays two roles in this paper:

  1. Exact termination locus: By Corollary 4.5, the Pellis expansion of x $\in$ Z[$\varphi$] terminates
  in finitely many steps with zero remainder. Thus Z[$\varphi$] is precisely the set of real
  numbers with finite Pellis expansions --- it is the fixed-point set of the expansion
  operator.

  2. Discrete state space for symbolic dynamics: The symbolic dynamical system of
  Section 6 is defined over the space X = T $\times$ ZN; the Lucas ring furnishes the discrete
  skeleton of X by supplying exactly-representable anchor points.

  The connection between Z[$\varphi$] and the structure of Penrose tilings / quasicrystals is not
  accidental: the inflation and deflation rules of Penrose tilings scale tile lengths by $\varphi$,
  and the number of tiles after each generation satisfies a Fibonacci recurrence --- the
  same arithmetic that governs termination of the Pellis expansion.

        Proposition 2.4 (Lucas closure under Trinity product). Z[$\varphi$] cap T is closed
        under the Trinity product (Definition 3.3). Proof. Any product T(lambda1) $\cdot
        T(lambda2) with both factors in Z[$\varphi$] remains in Z[$\varphi$] by ring closure; the
        monomial form of the product is confirmed directly. □

  This proposition corresponds to C-02 (Lucas closure) in the PhD program's Coq citation
  map, with Qed status as of the CLARA Provenance Audit 2026-05-12 (Appendix E).

  2.6 Complexity of a Representation
  The bit-complexity (or description length) of a Trinity label lambda = (n,k,p,m,q) is

  ell(lambda) = lceil log2(n+1) rceil + lceil log2(|k|+1) rceil + 1 + lceil log2(|p|+1) rceil
  + 1 + lceil log2(|m|+1) rceil + 1 + lceil log2(|q|+1) rceil + 1,

  where the +1 terms account for sign bits. The Pellis depth D of an expansion
  contributes approximately D log2($\varphi$-1) approx 0.694 D bits per term in the coefficient
  sequence (see Section 7). This definition follows the two-part coding model of Barron,
  Rissanen, and Yu (1998).

  3. Multiplicative Grammar of Trinity Monomials

\subsection{References}

- Vasilev, D. \textit{Flos Aureus Monograph: Trinity Constants}. DOI \href{https://zenodo.org/doi/10.5281/zenodo.19227877}{10.5281/zenodo.19227877}.
- CODATA 2018 recommended values: NIST \href{https://physics.nist.gov/cuu/Constants/}{https://physics.nist.gov/cuu/Constants/}.
- Heyrovska, R. Golden-angle FSC publications --- Heyrovský Institute of Physical Chemistry.
